
\begin{table*}[ht]
\centering
\caption{Country structural profiles for the 2007 baseline}
\label{tab:country_profiles_2007}
\scriptsize
\renewcommand{\arraystretch}{1.12}
\setlength{\tabcolsep}{4pt}
\begin{tabular}{lccccc lccccc}
\toprule
\textbf{Country} & \textbf{B1} & \textbf{B2} & \textbf{B3} & \textbf{B4} & \textbf{B5}
& \textbf{Country} & \textbf{B1} & \textbf{B2} & \textbf{B3} & \textbf{B4} & \textbf{B5} \\
\midrule
AUT & 2 & 4 & 5 & 3 & 2 & IRL & 4 & 4 & 2 & 4 & 1 \\
BEL & 1 & 2 & 4 & 4 & 1 & ITA & 1 & 3 & 2 & 1 & 2 \\
CAN & 3 & 3 & 4 & 5 & 5 & LTU & 5 & 5 & 1 & 3 & 3 \\
CZE & 4 & 3 & 3 & 1 & 5 & LUX & 5 & 5 & 3 & 3 & 1 \\
DEU & 2 & 1 & 4 & 2 & 3 & LVA & 5 & 3 & 1 & 2 & 2 \\
DNK & 4 & 5 & 5 & 4 & 5 & NLD & 3 & 5 & 4 & 4 & 4 \\
ESP & 4 & 1 & 3 & 3 & 1 & POL & 3 & 1 & 1 & 2 & 4 \\
EST & 5 & 5 & 2 & 5 & 5 & PRT & 1 & 2 & 2 & 1 & 1 \\
FIN & 3 & 2 & 5 & 5 & 3 & SVK & 4 & 1 & 1 & 1 & 2 \\
FRA & 2 & 2 & 4 & 3 & 3 & SVN & 5 & 4 & 3 & 2 & 4 \\
GBR & 1 & 4 & 3 & 5 & 5 & SWE & 3 & 3 & 5 & 4 & 4 \\
GRC & 1 & 1 & 1 & 2 & 2 & USA & 2 & 4 & 5 & 5 & 4 \\
HUN & 2 & 2 & 2 & 1 & 3 &  &  &  &  &  &  \\
\bottomrule
\end{tabular}
\vspace{2pt}
\begin{flushleft}
\scriptsize \textit{Notes:} B1 = debt capacity; B2 = employment strength; B3 = R\&D intensity; B4 = tertiary human capital; B5 = energy security. All variables are coded on a 1--5 ordinal scale, where higher values indicate a more favourable structural condition.
\end{flushleft}
\end{table*}
